% Pursuing Stacks (cotes 134-2 to 134-8): archive cote 134-4, site page 2 (PDF page 3), typescript page 90.
% typescript. Draft compared: ps3.tex:46-112.
% AI-transcribed by Claude Opus 5.5 (claude-opus-5-5, adaptive thinking, effort high), Message Batches API, batch ps1 (msgbatch_01Rv5hCaNTtUt4JHx1yqoHPL), 2026-11-12; not yet checked by eye.
\scanpage{134-4/2}
\mspage{(90)}

Test categories seem to play a similar role here as \inserted[typed]{(the spectra of)} discrete valuati-%
on rings in algebraic geometry -  they can be mapped into anywhere\hand{,}to
``test'' what is going on there  -  here it means,  they can be sent into
any modelizer (M,W)  (at least among the ones which we feel are the most
interesting) ,  by \xout{functors} ``test functors''  $i\colon A\longrightarrow M$  giving rise to
a model-preserving functor  $i^{*}$ :  $M\longrightarrow A\hat{}$,  allowing comparison of M with
an elementary modelizer $A\hat{}$.  As for the all-encompassing basic modelizer
(Cat),  it seems to play the opposite role in a sense,  at least with
respect to elementary modelizers\xout{\unsure{,}} $A\hat{}$,  which all admit modelizing maps
$i_A$ :  $A\hat{}\longrightarrow$(Cat).  As a matter of fact,  for given \over{\unsure{t}}{t}est category A i.e.
a given elementary modelizer $A\hat{}$,  I see for the time being just \uline{one} way
to get a modelizing functor from it to (Cat),  namely just the canonical
$i_A$ .  There is also a striking difference between the exactness properties
of the functors
\[ \text{(1)}\qquad i^{*} : M\longrightarrow A\hat{} \]
one way,  which commute to inverse limits,  and the functors
\[ \text{(2)}\qquad i_A\colon A\hat{}\longrightarrow\text{(Cat)} \]
in the other direction,  commuting to direct limits.  Another difference
is that we should not expect that the left adjoint  $i_!$ to $i^{*}$ be model pre-%
serving too (with the exception of the very special case when M=(Cat) and
$i\colon A\longrightarrow M$ = (Cat)  is the canonical functor $i_A$ ,  which appears as highly
non typical in this respect ),  whereas the right adjoint  $j_A=i_A^{\,*}$  of $i_A$
is model preserving,  thus $i_A$ is part of a pair $(i_A,j_A)$ of model preser-%
ving adjoint functors.

Of course,  we may want to compare\inserted{directly} \deleted{an} arbitrary modelizer M to (Cat)
by sending it into (Cat) by a modelizing functor  $M\longrightarrow$(Cat);  we get quite
naturally such a functor \hand{(for any given choice of test-functor $i\colon A\to M$)}
\[ \text{(3)}\qquad i_Ai^{*} :\ M\longrightarrow\text{(Cat)}\quad , \]
but this functor is not likely any more to commute \deleted{\unsure{to}} neither\inserted{to} direct nor\hand{ to}
inverse limits, even finite ones  -  and it isn\hand{'}t too clear that for a
modelizer M which is'nt elementary,  we have much chance to get a modeli-%
zing functor to (Cat) which is either left or right exact.  However,  the
functors\inserted{(3)} we've got,  whenever modelizing and if M is a \uline{strict} modelizer
(namely $M\longrightarrow W_M^{-1}M\simeq$(Hot)  commutes with finite products),  will commute
to finite products ``up to weak equivalence''.  Also the functors $i^{*}$ ,  al-%
though not right exact definitely,  have a tendency to commute to sums,\margin{also to fibered products $\times$}
and hence the same will hold (not only up to weak equivalence) for (3).

As for getting a modelizing functor  (Cat)$\longrightarrow$M ,  for a modelizer M

\margin{2}
